
\def\PUDcodigo{EME101}

\def\PUDcodacad{04506.52}

\def\PUDdisciplina{Corrosão e Proteção Anti-Corrosiva}

\def\PUDsemestre{Opt}

\def\PUDhoras{80}

%NOVOS ---------------------------------------------------
\def\PUDhorasTeoria{60}
\def\PUDhorasPratica{20}

\def\PUDinterECA{---}

\def\PUDatuaisECA{---}

\def\PUDrecursos{Projetor multimídia; Quadro branco e pincel.}

%----------------------------------------------------------

\def\PUDcreditos{4}

\def\PUDprerequisito{ \hyperref[PUD:EME108]{EME108} }

\def\PUDpreacad{ \hyperref[PUD:EME108]{04506.20} }

\def\PUDobjetivos{Entender a relação entre conceitos básicos de eletroquímica e os fenômenos responsáveis pela corrosão dos materiais metálicos;  Distinguir os fenômenos responsáveis pela corrosão dos materiais metálicos e os danos diretos ou indiretos causados à natureza pela corrosãRelacionar
as possíveis causas da corrosão;  Propor soluções para problemas de corrosão e seu impacto ambiental.}

\def\PUDementa{Importância Social e Econômica da Corrosão;  Conceitos Básicos de Eletroquímica Aplicada à
Corrosão;  Formas de Corrosão;  Corrosão Influenciada por Fatores Mecânicos;  Meios Corrosivos;  Heterogeneidades Responsáveis pela Corrosão Eletroquímica;  Passivação;  Oxidação a Altas Temperaturas;  Proteção Anti-Corrosiva.}

\def\PUDconteudo{ & Corrosão. 
\\ 
 & Oxidação-redução. 
\\ 
 & Pilhas eletroquímicas. 
\\ 
 & Formas de corrosão. 
\\ 
 & Mecanismos básicos de corrosão. 
\\ 
 & Meios corrosivos. 
\\ 
 & Heterogeneidades responsáveis por corrosão eletroquímica. 
\\ 
 & Corrosão galvânica. 
\\ 
 & Oxidação a altas temperaturas. 
\\ 
 & Corrosão associada a fatores mecânicos. 
\\ 
 & Tipos de proteção contra a corrosão. }

\def\PUDmetodo{Aulas expositórias que podem ser teóricas e/ou práticas, onde as práticas no laboratório serão marcadas no decorrer da disciplina.}

\def\PUDavalia{A avaliação será feita com aplicação de provas teóricas e/ou práticas, além da possibilidade de inclusão de trabalhos e seminários no decorrer da disciplina.}

\def\PUDbibbasica{ & \href{www.biblioteca.ifce.edu.br/index.asp?codigo_sophia=23938}{Gentil}, Vicente;  Corrosão.  5º ed.  Rio de Janeiro - RJ.  Editora LTC.  2007.  353p.  ISBN: 9788521615569.  
\\ 
 & \href{www.biblioteca.ifce.edu.br/index.asp?codigo_sophia=24027}{Ramanathan}, Lalgudi Venkataraman.  Corrosão e Seu Controle.  5º ed.  São Paulo - SP.  Editora HEMUS.  339p.  ISBN: 8528900010. 
\\ 
 & \href{www.biblioteca.ifce.edu.br/index.asp?codigo_sophia=44594}{Gemelli}, Enori.  Corrosão de Materiais Metálicos e Sua Caracterização.  Rio de Janeiro.  Editora LTC.  2001.  183p.  ISBN: 8521612907. }

\def\PUDbibcomplementar{ & \href{www.biblioteca.ifce.edu.br/index.asp?codigo_sophia=61962}{Pelliccione}, A.  S.;  Moraes, M.  F.;  Galvão, J.  L.  R.;  Mello, L.  A.;  Silva, E.  S.  Análise de Falhas em Equipamentos de Processo - Mecanismos de Danos e Casos Práticos.  2 ed.  Rio de Janeiro.  Editora Interciência.  2014.  386p.  ISBN: 9788571933286. 
\\ 
 & \href{www.biblioteca.ifce.edu.br/index.asp?codigo_sophia=62087}{Affonso}, Luiz Otávio Amaral.  Equipamentos Mecânicos – Análise de Falhas e Solução de Problemas.  3º ed.  Rio de Janeiro.  Editora Qualitymark.  2014.  372p.  ISBN: 9788541400367. 
\\ 
 & \href{www.biblioteca.ifce.edu.br/index.asp?codigo_sophia=62306}{Nunes}, Laerce de Paula.  Fundamentos de Resistência à Corrosão.  1ª edição.  Rio de Janeiro.  Editora  Interciência.  2007.  330p.  ISBN: 9788571931626.
\\
 &  \href{www.biblioteca.ifce.edu.br/index.asp?codigo_sophia=39995}{Chiaverini}, Vicente.  Aços e ferros fundidos: características gerais, tratamentos térmicos e principais tipos.  7º ed.  São Paulo, SP: ABM, 2012.  599p.  ISBN: 8586778486.
\\
 &  \href{www.	biblioteca.ifce.edu.br/index.asp?codigo_sophia=24421}{Callister Júnior}, William D.  Ciência e engenharia de materiais: uma introdução. 8º ed.  Rio de Janeiro, RJ: LTC, 2012.  817p.  ISBN: 9788521621249. }

\def\PUDautor{Francisco Nélio Costa Freitas}

\def\PUDdata{2014-05-21}

\def\PUDrevisao{2}

\def\PUDrevisor{Samuel Vieira Dias}

\def\PUDdatarevisao{2017-04-30}

\PUDcorpo
